\begin{figure}[!ht]
\centering
\scalebox{0.8}{ % Scale to 80%
\begin{tikzpicture}

% --- Part 1: C(0) ---
\begin{scope}
    % A set of 7 points (0-simplices)
    \node[dot] at (0.2, 0.8) {};
    \node[dot] at (1.0, 0.4) {};
    \node[dot] at (1.5, 1.5) {};
    \node[dot] at (2.1, 3.0) {};
    \node[dot] at (1.8, 3.8) {};
    \node[dot] at (0.5, 4.4) {};
    
    % Label for C(0)
    \node at (1.2, -0.8) {$\check{C}_0$};
\end{scope}

% --- Subset Symbol 1 ---
% \node at (4, 2.5) {$\subseteq$};

% --- Part 2: C(1) ---
\begin{scope}[xshift=5cm]
    \def\radius{0.65}
    
    % Define points as named nodes
    \coordinate (p1) at (0.2, 0.8); % Bottom component, point 1
    \coordinate (p2) at (1.0, 0.4); % Bottom component, point 2 (center)
    \coordinate (p3) at (1.5, 1.5); % Bottom component, point 3
    \coordinate (p4) at (2.1, 3.0); % Middle component, point 1
    \coordinate (p5) at (1.8, 3.8); % Middle component, point 2
    \coordinate (p6) at (0.5, 4.4); % Top isolated point

    % Draw circles around each point
    \foreach \i in {1,...,6} {
        \fill[lightgray, opacity=0.5] (p\i) circle (\radius);
    }

    \node[dot] at (p1) {};
    \node[dot] at (p2) {};
    \node[dot] at (p3) {};
    \node[dot] at (p4) {};
    \node[dot] at (p5) {};
    \node[dot] at (p6) {};

    % Draw edges (1-simplices)
    \draw (p1) -- (p2) -- (p3);
    \draw (p4) -- (p5);
    
    % Label for C(1)
    \node at (1.6, -0.8) {$\check{C}_1$};
\end{scope}

% --- Subset Symbol 2 ---
% \node at (9, 2.5) {$\subseteq$};

% --- Part 3: C(2) ---
\begin{scope}[xshift=10cm]
    \def\radius{0.9}
    
    % Define coordinates for the 7 points
    \coordinate (v1) at (0.2, 0.8) {}; % Bottom component, point 1
    \coordinate (v2) at (1.0, 0.4) {}; % Bottom component, point 2 (center)
    \coordinate (v3) at (1.5, 1.5) {}; % Bottom component, point 3
    \coordinate (c1) at (2.1, 3.0) {}; % Middle component, point 1
    \coordinate (c2) at (1.8, 3.8) {}; % Middle component, point 2
    \coordinate (c3) at (0.5, 4.4) {}; % Top isolated point
    
    % Fill the triangle (2-simplex) - drawn first to be in the background
    \fill[darkgray] (v1) -- (v2) -- (v3) -- cycle;
    
    % Draw circles around each point
    \foreach \p in {c1, c2, c3, v1, v2, v3} {
        \fill[lightgray, opacity=0.5] (\p) circle (\radius);
    }
    
    % Draw all edges (1-simplices)
    \draw (v3) -- (c1) -- (c2) -- (c3);
    \draw (v1) -- (v2) -- (v3) -- cycle;
    
    % Draw the points (0-simplices) - drawn last to be on top
    \foreach \p in {c1, c2, c3, v1, v2, v3} {
        \node[dot] at (\p) {};
    }
    
    % Label for C(2)
    \node at (1.8, -0.8) {$\check{C}_2$};
\end{scope}

\end{tikzpicture}
} % End scalebox
\caption{Esempio di un complesso di \v{C}ech, con gli $0$-simplessi (punti) in $\check{C}_0$, gli $1$-simplessi (archi) che si aggiungono in $\check{C}_1$, e un $2$-simplesso (triangolo) in $\check{C}_2$.}
\end{figure}